\chapter{General Context and Project Scope}
\label{ch:context}

\section{Introduction}

Modern information systems span servers, networks, applications, containers, cloud resources, and security tools. Each layer produces logs, metrics, and alerts. In a Security Operations Center (SOC), the challenge is to transform this data into understandable, prioritized, and actionable information.

The final-year internship project, titled \textbf{Intelligent Platform for Automated Monitoring and Security}, addresses this challenge through ARIA. It centralizes security and observability events, detects threats, explains incidents with an AI-assisted layer, and automates selected corrective actions with Ansible while preserving human validation for sensitive operations. This chapter defines ARIA's context: the host organization, the cybersecurity and SOC background, the existing situation and its limits, the problem statement, the objectives, the proposed solution, and the methodology and planning.

\section{Host Organization Presentation}

The internship was carried out at Huawei Technologies, the host organization for the ARIA project. Huawei is a telecommunications and technology company active in infrastructure, cloud, networks, and digital services. The administrative project sheet lists its Tunis office as Avenue de la Bourse, Immeuble Neifer, 1053 Les Berges du Lac, Tunis, Tunisia.

\begin{table}[H]
\centering
\caption{Administrative internship information from restitution documents}
\begin{tabularx}{\textwidth}{@{}p{0.28\textwidth}Y@{}}
\toprule
\textbf{Field} & \textbf{Value} \\
\midrule
Student & Ghazi Mabrouki \\
Institution / training & ESPRIT University -- ArcTic \\
Host organization & Huawei Technologies \\
Project context & ARIA Project / SOC, monitoring, cybersecurity, automation, AI \\
Internship period & January 12, 2026 -- July 12, 2026 \\
Company supervisors & Arfaoui Ibrahim and Zhong Yan \\
Academic supervisor & Mongi Dziri \\
\bottomrule
\end{tabularx}
\end{table}

The project sits at the intersection of SOC engineering, SOAR workflow design, backend development, observability, and infrastructure automation, covering cloud infrastructure, operational cybersecurity, log analysis, security supervision, and response automation.

\section{Cybersecurity and SOC Context}

A Security Operations Center is the team and tooling responsible for detecting, analyzing, and responding to security events across an organization's information system. It continuously receives telemetry from host-based detection agents, network intrusion detection engines, runtime security sensors, system metrics, application logs, and manual reports. Its value lies in turning this flood of raw signals into a small number of well-understood, actionable incidents.

Two operational disciplines are central to this project. \textbf{Monitoring and detection} provides visibility over the host, network, runtime, and resource layers. \textbf{SOAR} (Security Orchestration, Automation, and Response) sits above detection and coordinates investigation and remediation, automating repetitive analysis and corrective actions to reduce manual effort and response time while keeping the human analyst responsible for risky decisions. ARIA operates in this SOAR domain: it consumes monitoring output and adds the orchestration, intelligence, and controlled automation layer a SOC needs.

The internship work was organized around three complementary technical fronts that together form a complete detection-to-response chain:

\begin{itemize}[leftmargin=1cm]
\item \textbf{A security monitoring foundation:} Elasticsearch, Kibana, Wazuh, Suricata, Falco, Filebeat, Telegraf, and Redis, providing visibility over the host, network, runtime, and resource layers.
\item \textbf{A cloud infrastructure:} a Huawei Cloud environment hosting the monitoring stack and the ARIA platform in a controlled, isolated environment.
\item \textbf{A backend intelligence layer:} the ARIA platform that transforms raw Elasticsearch documents into normalized alerts, incidents, AI-assisted investigations, controlled remediation, verification, archives, and operator actions.
\end{itemize}

The complete high-level chain can be summarized as follows:

\begin{center}
\texttt{Security tools -> Elasticsearch -> Backend intelligence -> Frontend SOC operations}
\end{center}

\section{State of the Art and Identified Gap}

Several commercial and open-source platforms address parts of the SOC workflow. Table~\ref{tab:state-of-the-art} summarizes the most relevant ones.

\begin{table}[H]
\centering
\caption{State-of-the-art SOC and SOAR solutions}
\label{tab:state-of-the-art}
\begin{tabularx}{\textwidth}{@{}p{0.22\textwidth}p{0.34\textwidth}Y@{}}
\toprule
\textbf{Solution} & \textbf{Strengths} & \textbf{Limitations} \\
\midrule
Splunk ES & Powerful SIEM, rich app ecosystem & High cost, heavy infrastructure \\
IBM QRadar & Correlation, asset database & Complex deployment and tuning \\
Microsoft Sentinel & Cloud-native, built-in AI alerts & Vendor lock-in, cloud dependency \\
Palo Alto XSOAR & Large playbook marketplace & Expensive, steep learning curve \\
Wazuh + ELK & Open source, good visibility & No built-in SOAR or AI response layer \\
Falcon / SentinelOne & Strong EDR capabilities & EDR-centric, limited cross-layer correlation \\
\bottomrule
\end{tabularx}
\end{table}

None combine multi-layer monitoring (host, network, runtime, and metrics), alert enrichment, incident correlation, AI-generated playbooks, and human-controlled Ansible remediation in a single open, local-deployment platform. This gap motivates ARIA.

\begin{center}
\textit{Identified gap: there is no open, integrated solution that turns heterogeneous security and observability telemetry into actionable, approved, and verifiable incidents.}
\end{center}

\section{Study of the Existing Situation}

Before ARIA, the environment relied on several specialized tools: Wazuh for host security visibility, Suricata for suspicious network behavior, Falco for runtime behavior, Telegraf for metrics, Filebeat for log forwarding, and Elasticsearch/Kibana for storage and visualization. These tools are powerful individually, but a SOC workflow requires more than isolated visibility.

\begin{table}[H]
\centering
\caption{Existing sources and their operational role}
\begin{tabularx}{\textwidth}{@{}p{0.20\textwidth}Y@{}}
\toprule
\textbf{Source / tool} & \textbf{Role in the environment} \\
\midrule
Wazuh & Host and endpoint security monitoring, agent alerts, authentication and system activity. \\
Suricata & Network IDS/IPS visibility through EVE JSON events and protocol/security signatures. \\
Falco & Runtime and container-oriented detection based on system calls and abnormal behavior. \\
Telegraf & Infrastructure metrics such as CPU, memory, disk, network, load, processes, and sockets. \\
Filebeat / Falcosidekick & Collection and forwarding layer toward Elasticsearch. \\
Elasticsearch / Kibana & Central indexing, search, visualization, and dashboarding foundation. \\
\bottomrule
\end{tabularx}
\end{table}

\diagramimage{phase1_visibility_layers.png}{Monitoring visibility layers and tool roles}

The mission was not to replace the monitoring tools but to fill the operational layer above them. A useful comparison is therefore between three approaches: individual monitoring consoles, a search and visualization stack, and a SOAR-style response platform.

\begin{table}[H]
\centering
\caption{Comparison of existing approaches and ARIA}
\begin{tabularx}{\textwidth}{@{}p{0.22\textwidth}p{0.22\textwidth}p{0.22\textwidth}Y@{}}
\toprule
\textbf{Approach} & \textbf{Strengths} & \textbf{Limits} & \textbf{ARIA contribution} \\
\midrule
Tool-by-tool consoles & Each tool keeps its native detail and configuration model. & Analysts must manually switch context between Wazuh, Suricata, Falco, metrics, and logs. & Centralizes normalized operational views while preserving source evidence. \\
Elasticsearch and Kibana dashboards & Powerful search, visualization, and index exploration. & Raw events remain heterogeneous and response workflow is manual. & Adds normalization, incidents, investigations, AI assistance, approval, execution, verification, and archive. \\
Manual SOC runbooks & Human judgment remains high and risky actions are controlled. & Response time depends on expert availability and evidence collection is repetitive. & Automates repetitive analysis and playbook preparation while keeping approval and audit controls. \\
\bottomrule
\end{tabularx}
\end{table}

\section{Limits of the Existing Situation}

The initial environment produced useful telemetry, but had several operational limitations:

\begin{itemize}[leftmargin=1cm]
\item security and supervision tools produced heterogeneous data formats;
\item alerts were numerous and sometimes noisy, creating alert fatigue;
\item raw technical alerts were difficult to interpret quickly;
\item correlation between host, network, runtime, and metric events required manual effort;
\item incident explanation depended heavily on expert availability;
\item corrective actions were often manual and not consistently traceable;
\item the complete lifecycle from alert to remediation was dispersed across tools.
\end{itemize}

These limits share a single root cause: the tools deliver excellent \emph{visibility} but no shared \emph{operational language} and no \emph{controlled response} layer. Each tool uses its own format and stops at detection, leaving the analyst to connect, interpret, decide, and verify everything by hand.

\section{Problem Statement}

The problem can be formulated as follows: how can a centralized platform collect heterogeneous security and monitoring data, analyze and correlate events automatically, explain incidents through AI assistance, and automate corrective actions while keeping human control over sensitive decisions?

Several sub-problems follow from this question. Raw alert formats differ across sources, making correlation difficult. Repeated or noisy alerts can hide important events. The analyst needs context such as IP information, MITRE ATT\&CK techniques, source history, target host, and related incidents. Automatic remediation can be dangerous if it executes unvalidated commands or targets the wrong host. Finally, the result of a fix must be verified rather than assumed.

\section{Project Objectives}

To answer this problem, the project pursued the following objectives:

\begin{itemize}[leftmargin=1cm]
\item Build and stabilize a centralized monitoring foundation on cloud infrastructure so that real host, network, runtime, and resource telemetry reaches Elasticsearch reliably.
\item Implement a multi-source alert pipeline that polls Elasticsearch and normalizes alerts from security and observability sources into a common model.
\item Enrich and correlate alerts using GeoIP information, MITRE ATT\&CK mappings, Sigma noise rules, deduplication, and incident grouping.
\item Provide an AI-assisted response layer able to produce investigation summaries, narratives, risk assessments, and Ansible playbooks.
\item Add approval, execution, verification, archiving, audit, and safety mechanisms around remediation workflows.
\item Expose the platform through a FastAPI backend and a Next.js dashboard for SOC analysts and administrators.
\item Onboard new servers automatically as scoped monitored assets and prepare the solution for production deployment.
\item Validate the implementation through automated tests, compile/build checks, API checks, and runtime verification reports.
\end{itemize}

\section{Proposed Solution: ARIA}

ARIA is an intelligent SOC/SOAR platform that connects the monitoring foundation built in Chapter~2 to a backend intelligence layer and a Next.js dashboard. It performs four main functions:

\begin{itemize}[leftmargin=1cm]
\item \textbf{Ingest and normalize:} poll Elasticsearch indices and turn heterogeneous raw documents from Wazuh, Suricata, Falco, Telegraf, and generic sources into a common alert model.
\item \textbf{Enrich and correlate:} add GeoIP and MITRE ATT\&CK context, remove duplicates and noise, and group related alerts into incidents.
\item \textbf{Assist and control response:} generate AI-assisted investigation summaries, risk scores, narratives, and Ansible playbooks; require analyst approval and admin-secret validation before any sensitive action.
\item \textbf{Verify and archive:} execute playbooks, re-check evidence, archive closed cases, and persist a complete audit trail.
\end{itemize}

\diagramimage{phase3_global_backend.png}{ARIA backend intelligence layer}

This design keeps the analyst at the center of the decision process. AI accelerates understanding and proposes playbooks, while the platform controls execution through approval, admin-secret protection, safety checks, Ansible execution records, fix verification, and audit persistence.

The solution was built in three incremental phases. Phase~1 established the security monitoring foundation. Phase~2 provisioned the Huawei Cloud infrastructure that hosts it. Phase~3 implemented the ARIA backend intelligence and dashboard on top of the telemetry, then extended it with multi-server onboarding and production-oriented hardening.

\diagramimage{phase1_global_architecture.png}{Global architecture and data flow of the solution}

From the user's point of view, the scope of ARIA is best expressed as a global use-case diagram (Figure~\ref{fig:global-usecase}). It shows the primary human actors---the SOC analyst and the administrator---and the system actors ARIA coordinates: monitored servers, Elasticsearch, the AI provider, the Ansible executor, and the monitoring tools. The use cases fall into monitoring and triage, investigation and remediation, AI interaction, and administration, the four areas detailed in Chapters~\ref{ch:context} to~6. The diagram defines the functional boundary of the platform: everything inside is an ARIA responsibility; every external actor is an integration, not a replacement.

\umlfigure{global_use_case.png}{Global use-case diagram of the ARIA platform}{fig:global-usecase}

\section{Methodology and Planning}

\subsection{Choice and Justification}

We selected Scrum as the guiding methodology, adapted for a single-developer internship context. The justification rests on three factors:

\begin{enumerate}[leftmargin=1cm]
  \item \textbf{Emergent requirements.} At the start of the internship the precise behavior of the AI playbook engine, the optimal enrichment taxonomy (MITRE, GeoIP, Sigma, IOC, whitelist), and the multi-server scoping model were unknown. Scrum allowed us to refine these through working software reviewed at the end of each sprint.
  \item \textbf{Frequent supervisor feedback.} Bi-weekly sprint reviews with the company supervisors, plus checkpoint reviews with the academic supervisor, ensured alignment with Huawei's production constraints and ESPRIT's academic expectations and prevented divergence from operational reality.
  \item \textbf{Incremental validation.} Each sprint produced a demonstrable increment: event detection and normalization in Sprint~5, incident correlation in Sprint~6, AI-assisted investigations in Sprint~7, controlled remediation loops in Sprint~9, and multi-server onboarding in Sprint~11. This made integration issues visible early and gave continuous evidence that the objectives were achievable.
\end{enumerate}

The Scrum Guide \cite{scrum_guide} provides the theoretical basis; the ceremonies were scaled down to fit a six-month internship and a single main developer.

\subsection{Adapted Scrum Process}

The internship was organized into thirteen two-week sprints from January~12 to July~12, 2026. The process kept the essential Scrum artifacts and events but removed overhead that would be excessive for one intern:

\begin{itemize}[leftmargin=1cm]
  \item \textbf{Product backlog:} a single prioritized list of features, described in Table~\ref{tab:backlog}.
  \item \textbf{Sprint backlog:} at the start of each sprint a small subset of product-backlog items was selected, decomposed into concrete tasks, and written in the internship logbook.
  \item \textbf{Daily tracking:} the logbook recorded daily progress, blockers, and decisions instead of a formal daily stand-up.
  \item \textbf{Sprint review:} every two weeks the current increment was demonstrated to the company supervisors; feedback was captured as new or refined backlog items.
  \item \textbf{Sprint retrospective:} a short internal review at the end of each sprint identified what to keep, adjust, or stop for the next sprint.
\end{itemize}

Documentation, validation, and report writing were not postponed to the end; they ran continuously, so each sprint closed with both working code and updated evidence.

\subsection{Sprint Planning}

Table~\ref{tab:sprints} lists the thirteen sprints, their dates, objectives, deliverables, and the report chapters they feed.

{\small\setlength{\tabcolsep}{2pt}
\begin{longtable}{@{}p{0.07\textwidth}@{}p{0.12\textwidth}@{}p{0.16\textwidth}@{}p{0.46\textwidth}@{}p{0.10\textwidth}@{}}
\caption{Sprint planning (January 12 -- July 12, 2026)}\label{tab:sprints}\\
\toprule
\textbf{Sprint} & \textbf{Dates} & \textbf{Objective} & \textbf{Key deliverables} & \textbf{Chapters} \\
\midrule
\endfirsthead
\toprule
\textbf{Sprint} & \textbf{Dates} & \textbf{Objective} & \textbf{Key deliverables} & \textbf{Chapters} \\
\midrule
\endhead
S1 & 12 Jan -- 25 Jan & Research, context, requirements & Host organization study, problem statement, objectives, methodology & 1 \\
S2 & 26 Jan -- 08 Feb & Cloud environment setup & VPC, subnet, security group, EIP, ECS host, storage & 2 \\
S3 & 09 Feb -- 22 Feb & Monitoring foundation & Elasticsearch, Kibana, Wazuh, Suricata, Falco, Filebeat, Telegraf, Redis & 2 \\
S4 & 23 Feb -- 08 Mar & Architecture \& design & ARIA backend skeleton, domain model, API design, architecture diagrams & 3 \\
S5 & 09 Mar -- 22 Mar & Alert ingestion & Cursor-based polling, source mapping, normalized alert persistence & 3 \\
S6 & 23 Mar -- 05 Apr & Enrichment \& correlation & GeoIP, MITRE, Sigma, IOC, whitelist, incident grouping & 3 \\
S7 & 06 Apr -- 19 Apr & AI response & AI summary, narrative, risk assessment, playbook generation & 3/4 \\
S8 & 20 Apr -- 03 May & Dashboard \& analyst workflow & Alerts, incidents, investigations UI & 4 \\
S9 & 04 May -- 17 May & Safe remediation & Approval workflow, admin secrets, Ansible execution, verification & 4 \\
S10 & 18 May -- 31 May & Runtime \& infrastructure & Runtime diagnostics, metrics anomalies, IPS visualization & 4 \\
S11 & 01 Jun -- 14 Jun & Multi-server onboarding & Asset registry, bootstrap, asset scoping & 5 \\
S12 & 15 Jun -- 28 Jun & Validation \& testing & Build checks, API tests, live evidence, automated tests & 6 \\
S13 & 29 Jun -- 12 Jul & Documentation \& final report & Report corrections, defense preparation, final submission & 1/6/Concl. \\
\bottomrule
\end{longtable}
\par}

\subsection{Product Backlog}

The product backlog in Table~\ref{tab:backlog} summarizes the features prioritized and implemented across the sprints. Each item is mapped to the sprint(s) in which it was delivered and to the acceptance criteria used to confirm it.

{\small\setlength{\tabcolsep}{2pt}
\begin{longtable}{@{}p{0.06\textwidth}@{}p{0.16\textwidth}@{}p{0.48\textwidth}@{}p{0.09\textwidth}@{}p{0.10\textwidth}@{}}
\caption{Product backlog with sprint mapping and acceptance criteria}\label{tab:backlog}\\
\toprule
\textbf{ID} & \textbf{Feature} & \textbf{Description / Acceptance criteria} & \textbf{Priority} & \textbf{Sprint(s)} \\
\midrule
\endfirsthead
\toprule
\textbf{ID} & \textbf{Feature} & \textbf{Description / Acceptance criteria} & \textbf{Priority} & \textbf{Sprint(s)} \\
\midrule
\endhead
PB-01 & Monitoring foundation & Install and stabilize the central security and observability stack on cloud infrastructure. \textbf{Acceptance:} all services reachable; Kibana shows indices; health checks pass. & High & S2--S3 \\
PB-02 & Alert ingestion & Poll Elasticsearch sources and persist normalized alerts. \textbf{Acceptance:} cursor advances; no duplicate processing; normalized alerts stored in SQLite. & High & S5 \\
PB-03 & Alert enrichment & Add GeoIP, MITRE, Sigma, severity, IOC, and whitelist logic. \textbf{Acceptance:} enrichment fields populated; whitelist suppresses false positives. & High & S6 \\
PB-04 & Incident correlation & Group related alerts into incidents and investigations. \textbf{Acceptance:} correlation rules configurable; related alerts appear under one incident. & High & S6 \\
PB-05 & AI response & Generate summary, narrative, risk assessment, and playbook. \textbf{Acceptance:} investigation contains all fields; malformed LLM output handled gracefully. & High & S7 \\
PB-06 & Analyst workflow & Approve, decline, edit, execute, archive, and audit investigations. \textbf{Acceptance:} state machine enforces allowed transitions; audit records persisted. & High & S8--S9 \\
PB-07 & Safe remediation & Validate playbooks, protect dangerous actions, execute via Ansible, verify fixes. \textbf{Acceptance:} admin-secret gate; dry-run available; execution log and verification saved. & High & S9 \\
PB-08 & Dashboard & Provide alerts, incidents, investigations, assistant, IPS, metrics, and settings pages. \textbf{Acceptance:} all pages render; API checks pass; production build succeeds. & High & S8, S10 \\
PB-09 & Runtime remediation & Diagnose runtime findings and staged remediation actions. \textbf{Acceptance:} observe cases never produce corrective playbooks; remediation requires explicit approval. & Medium & S10 \\
PB-10 & Infrastructure monitoring & Poll Telegraf metrics and detect anomalies. \textbf{Acceptance:} metrics ingested; anomaly rules trigger investigations. & Medium & S10 \\
PB-11 & Multi-server onboarding & Add asset registry, bootstrap monitored servers, and scope data by server. \textbf{Acceptance:} new asset added; bootstrap installs required services; endpoints filter by \codepath{asset\_id}. & Medium & S11 \\
PB-12 & Validation & Compile, test, build, and document verification results. \textbf{Acceptance:} automated tests pass; build checks green; evidence screenshots collected. & High & S12--S13 \\
\bottomrule
\end{longtable}
\par}

The complete planning is visualized in the Gantt chart of Figure~\ref{fig:gantt}.

\clearpage
\begin{landscape}
\thispagestyle{fancy}
\begin{figure}[p]
\centering
\newcommand{\ganttleft}{-5.5}
\newcommand{\ganttwidth}{26}
\newcommand{\gantttop}{15}
\newcommand{\rowzero}{13.5}

\begin{tikzpicture}[x=0.62cm,y=0.78cm, >=stealth]
  % ---- Custom palette ----
  \definecolor{lavender}{RGB}{230,230,250}
  \definecolor{paleyellow}{RGB}{255,250,205}
  \definecolor{palegray}{RGB}{240,240,240}

  % ---- Styles ----
  \tikzset{
    gridline/.style={gray!35, very thin},
    sprintline/.style={gray!65, thin},
    milestone/.style={diamond, fill=violet!70, draw=violet!50!black, line width=0.4pt, inner sep=2.2pt},
    milestonelabel/.style={font=\tiny\bfseries, text=violet!80!black, anchor=south},
    taskbar/.style={rounded corners=2pt, line width=0.35pt, draw=gray!55!black},
    deliverbar/.style={rounded corners=2pt, line width=0.35pt, draw=gray!55!black},
  }

  % ---- Title ----
  \node[font=\bfseries\Large] at ({\ganttwidth/2},\gantttop+1.3)
    {Adaptive Response Intelligence Automation --- Project Plan};

  % ---- Month header ----
  \draw[sprintline] (0,\rowzero+0.55) -- (0,\rowzero-0.55);
  \node[font=\small\bfseries] at (1.43,\rowzero) {Jan};
  \draw[sprintline] (2.857,\rowzero+0.55) -- (2.857,\rowzero-0.55);
  \node[font=\small\bfseries] at (4.857,\rowzero) {Feb};
  \draw[sprintline] (6.857,\rowzero+0.55) -- (6.857,\rowzero-0.55);
  \node[font=\small\bfseries] at (9.071,\rowzero) {Mar};
  \draw[sprintline] (11.286,\rowzero+0.55) -- (11.286,\rowzero-0.55);
  \node[font=\small\bfseries] at (13.429,\rowzero) {Apr};
  \draw[sprintline] (15.571,\rowzero+0.55) -- (15.571,\rowzero-0.55);
  \node[font=\small\bfseries] at (17.786,\rowzero) {May};
  \draw[sprintline] (20.0,\rowzero+0.55) -- (20.0,\rowzero-0.55);
  \node[font=\small\bfseries] at (22.143,\rowzero) {Jun};
  \draw[sprintline] (24.286,\rowzero+0.55) -- (24.286,\rowzero-0.55);
  \node[font=\small\bfseries] at (25.143,\rowzero) {Jul};
  \draw[sprintline] (\ganttwidth,\rowzero+0.55) -- (\ganttwidth,\rowzero-0.55);

  % ---- Row helper ----
  % #1 = row index (0 at top), #2 = label, #3 = start week, #4 = end week,
  % #5 = bar fill, #6 = bar label, #7 = band color
  \newcommand{\ganttrow}[7]{%
    \fill[#7, opacity=0.22] (\ganttleft,{\rowzero-#1-0.48}) rectangle (\ganttwidth,{\rowzero-#1+0.48});
    \draw[gridline] (\ganttleft,{\rowzero-#1-0.5}) -- (\ganttwidth,{\rowzero-#1-0.5});
    \node[anchor=east, font=\small, text width=5cm, align=right] at (-0.3,{\rowzero-#1}) {#2};
    \fill[#5, taskbar] (#3,{\rowzero-#1-0.28}) rectangle (#4,{\rowzero-#1+0.28});
    \node[anchor=west, font=\tiny\bfseries, text=white] at (#3+0.12,{\rowzero-#1}) {#6};
  }

  % ---- Rows ----
  \ganttrow{0}{Research / Discovery}{0}{2}{gray!55}{Context research}{lavender!40}
  \ganttrow{1}{Requirements \& Specification}{0}{3}{gray!55}{Requirements}{yellow!30}
  \ganttrow{2}{Cloud \& Linux Environment}{2}{5}{gray!55}{Cloud setup}{gray!20}
  \ganttrow{3}{Monitoring Foundation}{4}{7}{gray!55}{Monitoring stack}{lavender!40}
  \ganttrow{4}{Architecture \& Design}{6}{9}{blue!45!violet!65}{Architecture}{yellow!30}
  \ganttrow{5}{Backend: Ingestion}{8}{11}{gray!55}{Ingestion}{gray!20}
  \ganttrow{6}{Backend: Enrichment \& Correlation}{10}{13}{gray!55}{Enrichment}{lavender!40}
  \ganttrow{7}{Backend: AI Response}{12}{15}{gray!55}{AI response}{yellow!30}
  \ganttrow{8}{Dashboard \& SOC Workflows}{14}{18}{gray!55}{Dashboard}{gray!20}
  \ganttrow{9}{Safe Remediation}{16}{19}{gray!55}{Remediation}{lavender!40}
  \ganttrow{10}{Runtime \& Infrastructure Monitoring}{18}{21}{gray!55}{Runtime / Infra}{yellow!30}
  \ganttrow{11}{Multi-server Onboarding \& Hardening}{20}{23}{gray!55}{Multi-server}{gray!20}
  \ganttrow{12}{Testing \& Validation}{22}{25}{blue!45!violet!65}{Validation}{lavender!40}
  \ganttrow{13}{Documentation \& Final Report}{23}{26}{blue!45!violet!65}{Documentation}{yellow!30}

  % ---- Bottom separator ----
  \draw[gridline] (\ganttleft,{\rowzero-14+0.5}) -- (\ganttwidth,{\rowzero-14+0.5});

  % ---- Vertical grid: every week thin, every sprint thicker ----
  \foreach \w in {0,1,...,\ganttwidth}{%
    \draw[gridline] (\w,\rowzero-0.5) -- (\w,{\rowzero-14+0.5});
  }
  \foreach \w in {0,2,4,...,\ganttwidth}{%
    \draw[sprintline] (\w,\rowzero-0.5) -- (\w,{\rowzero-14+0.5});
  }

  % ---- Sprint labels ----
  \foreach \s/\w in {1/0,2/2,3/4,4/6,5/8,6/10,7/12,8/14,9/16,10/18,11/20,12/22,13/24}{%
    \node[font=\tiny, below, text=gray!80!black] at (\w+1,{\rowzero-14+0.4}) {S\s};
  }

  % ---- Milestones ----
  % Architecture validated (end of architecture design)
  \node[milestone] (m1) at (9,\rowzero-4+0.55) {};
  \node[milestonelabel] at (m1.north) {Architecture validated};
  % PoC completed (first end-to-end alert-to-remediation loop)
  \node[milestone] (m2) at (13,\rowzero-7+0.55) {};
  \node[milestonelabel] at (m2.north) {PoC completed};
  % Test campaign completed
  \node[milestone] (m3) at (25,\rowzero-12+0.55) {};
  \node[milestonelabel] at (m3.north) {Test campaign completed};
  % Documentation completed
  \node[milestone] (m4) at (25.5,\rowzero-13+0.55) {};
  \node[milestonelabel, anchor=south east] at (m4.north west) {Documentation completed};
  % Final report submitted
  \node[milestone] (m5) at (26,\rowzero-13+0.55) {};
  \node[milestonelabel, anchor=south west] at (m5.north east) {Final report submitted};

  % ---- Current progress line (around 23 June 2026 = week 23.29) ----
  \draw[red, line width=1pt, dashed] (23.29,\rowzero+0.6) -- (23.29,{\rowzero-14+0.2});
  \node[font=\small\bfseries, text=red, anchor=south] at (23.29,\rowzero+0.65) {Current progress};

  % ---- Legend ----
  \fill[gray!55, taskbar] (\ganttleft,-1.1) rectangle (\ganttleft+2,-0.7);
  \node[anchor=west, font=\small] at (\ganttleft+2.2,-0.9) {Work package};

  \fill[blue!45!violet!65, taskbar] (\ganttleft+5.5,-1.1) rectangle (\ganttleft+7.5,-0.7);
  \node[anchor=west, font=\small] at (\ganttleft+7.7,-0.9) {Key deliverable};

  \node[milestone] at (\ganttleft+11.5,-0.9) {};
  \node[anchor=west, font=\small] at (\ganttleft+12,-0.9) {Milestone};

  \draw[red, line width=1pt, dashed] (\ganttleft+15,-0.9) -- (\ganttleft+17,-0.9);
  \node[anchor=west, font=\small, text=red] at (\ganttleft+17.2,-0.9) {Current progress};
\end{tikzpicture}
\caption{Project Gantt chart showing the full internship timeline across project phases and planned deliverables.}
\label{fig:gantt}
\end{figure}
\end{landscape}


\section{Conclusion}

This chapter presented the internship context, the host organization, the cybersecurity and SOC background, the initial monitoring environment and its limitations, the problem statement, the project objectives, the proposed ARIA solution, and the methodology and planning. The next chapter starts the engineering story at its foundation: the Huawei Cloud environment and the central monitoring stack that produces and stores the real telemetry ARIA consumes.
